\section{Why training compression is different}
\label{sec:introduction}

The primary question of this paper is simple:
\emph{how much can one compress the training dynamics of a neural network?}
The analogous question for a single trained snapshot is much easier to pose.
One may prune or factor its weights, distill its predictions, or approximate
its input--output map.  None of these operations alone supplies a compressed
state for continued training.  Two parameter states can represent the same
function on every input and still have different tangent features.  Their
next gradient updates can therefore differ.  A small one-step discrepancy can
then feed back into later activations, gradients, and updates.

We call a compression of training \emph{global on finite horizons} if, for
each finite $T$ and tolerance $\varepsilon$, one can choose a finite
compression order so that one autonomous compressed trajectory stays within
$\varepsilon$ of the dense trajectory for every $0\le t\le T$.  This is
stronger than approximating the vector field at one state and weaker than a
single error bound uniform over $0\le t<\infty$.  The state dimension should
remain constant as elapsed training time grows.  It may still depend on the
requested accuracy, dataset size, depth, and the representation class.

This question is scientifically useful even before it produces a smaller
implementation.  An accurate autonomous compression identifies information
that is sufficient for future learning, not merely for current prediction.
It can expose which aspects of the dense trajectory carry memory and which
details may be discarded.  Our nonnegotiable target retains three ingredients
together: more than one hidden layer, nonlinear responses that remain active,
and feature learning rather than a fixed-kernel linearization.

Wide-network limits offer a natural setting for this program.  With an
appropriate maximal-update or mean-field scaling, nontrivial feature motion
survives at large width \citep{yang2021tensor,nguyen2023rigorous}.  Empirical
averages may converge to deterministic fields, much as particle systems lead
to continuum descriptions.  Existing tensor-program, mean-field, and
dynamical-field theories provide powerful exact descriptions of such limits
\citep{yang2021tensor,nguyen2023rigorous,bordelon2022selfconsistent}.  In the
deep feature-learning settings closest to ours, however, exact descriptions
retain operators, two-time kernels, or a growing family of historical
responses.  They do not by themselves show that the learned dynamics can be
approximated by a fixed amount of history state throughout training.

The observation behind our construction is elementary.  For every internal
layer, gradient flow writes the learned matrix increment as an integral of
rank-one products,
\begin{equation}
 W_\ell(t)-W_{\ell,0}
 =-\frac{2}{nm}\sum_{a=1}^{m}\int_0^t
 r_a(s)\,\delta_{\ell,a}(s)h_{\ell-1,a}(s)^\top\,\od s .
 \label{eq:intro-history}
\end{equation}
Exact history storage grows with time.  We instead project the forward and
backward response paths onto the first $P$ shifted Legendre modes and evolve
their coefficients online.  The closure reconstructs the learned interaction
from these moment pairs.  Shifted Legendre memory is closely related to the
online polynomial-projection principle developed by HiPPO
\citep{gu2020hippo}.  The contribution here is to couple such memory to the
forward/backward feedback of deep gradient flow, preserve the reused
initialized operator in both orientations, and prove trajectory tracking for
the resulting autonomous nonlinear closure.

Two clocks produce two guarantees.  The first advances at residual RMS.  It
removes the residual factor from forward response speed and gives an
$O(P^{-1})$ state bound at arbitrary fixed depth.  The second advances at
residual activity plus the arclength of all stored forward and backward
responses.  It also keeps the original learning measure inside the history
integral.  Both histories then have first-order projection error, and their
outer-product error is $O(P^{-2})$.

The main contributions are as follows.
\begin{enumerate}[leftmargin=*,itemsep=0.35em]
\item We give explicit autonomous response-memory ODEs for a fixed number of
      moments per neuron, sample, and internal layer.  They retain each
      initialized matrix $W_{\ell,0}$ and its transpose exactly and never use
      a dense reference trajectory.
\item For arbitrary fixed width, finite dataset, fixed depth, arbitrary finite
      initialization, and any prescribed finite horizon, we prove uniform
      dense-trajectory error $O(P^{-1})$ for the activity clock and
      $O(P^{-2})$ for the response clock.  The proof controls the closure's
      own nonlinear feedback by an exact accumulator identity and a first-exit
      argument.
\item We distinguish three representation classes---scalar, population, and
      operator dynamics---and record a restricted exact scalar nonclosure
      theorem.  The response closure compresses the learned state to
      population fields but remains a hybrid with fixed initialized operators.
\item We give complete state, rank, and per-step cost accounting.  The learned
      correction has $O(nmP)$ moving coordinates per internal layer, while the
      exact initialized operators still require $O(n^2)$ storage and dense
      actions.  No total-model or generic runtime improvement is claimed.
\item We test the construction on sparse circle tasks, deeper networks,
      several activations, sphere tasks, and MNIST.  The experiments show high
      fidelity at small order in many cases, together with order reversals,
      numerical sensitivity, and difficult cases where low order fails.
\end{enumerate}

The distinction between proved and prospective population claims matters.
The central approximation theorem in this paper is pathwise at fixed finite
width.  Combining it with a separately established population limit gives a
diagonal sequence of closure orders converging to the same population
observables.  We do not prove that one fixed order works uniformly in width,
or that the closure constants remain controlled for all training time.
